% Generated draft: requires mathematical and source review.
\documentclass[11pt]{article}
\usepackage[margin=25mm]{geometry}
\usepackage{fontspec}
\setmainfont{Latin Modern Roman}
\usepackage{amsmath,amssymb,mathtools,unicode-math}
\setmathfont{Latin Modern Math}
\usepackage{enumitem,needspace,xurl,hyperref}
\setlength{\parindent}{0pt}
\setlength{\parskip}{5pt}
\setlength{\emergencystretch}{3em}
\newcommand{\N}{\mathbb{N}}
\newcommand{\Z}{\mathbb{Z}}
\newcommand{\Q}{\mathbb{Q}}
\newcommand{\R}{\mathbb{R}}
\newcommand{\C}{\mathbb{C}}
\newcommand{\F}{\mathbb{F}}
\newcommand{\A}{\mathbb{A}}
\newcommand{\PP}{\mathbb{P}}
\newcommand{\E}{\mathbb{E}}
\newcommand{\eps}{\varepsilon}
\newcommand{\dd}{\,\mathrm{d}}
\begin{document}
\section*{UnsolvedMath ID: 800426}
\textbf{Source problem code:} catalogue entry
\subsection{Definitions and mathematical statement}
For a polynomial \(f\in\Z[x]\), define its iterates by \(f^{\circ 0}(x)=x\) and \(f^{\circ(n+1)}=f\circ f^{\circ n}\) for integers \(n\geq 0\). The forward orbit of \(\alpha\in\Z\) is
\[
\mathcal O_f(\alpha)=\{f^{\circ n}(\alpha):n\in\Z_{\geq 0}\}.
\]
The point \(\alpha\) is called wandering if this orbit is infinite. Let \(\mathcal P=\{p\in\Z:p>1\text{ and }p\text{ is prime}\}\), using the usual convention of positive rational primes.

\textbf{Conjecture 49 (supplied source).} For every polynomial \(f\in\Z[x]\) with \(\deg(f)\geq 2\) and every wandering point \(\alpha\in\Z\) for \(f\),
\[
\#\bigl(\mathcal O_f(\alpha)\cap\mathcal P\bigr)<\infty.
\]
The local registry labels this conjecture open; that label is supplied source metadata, not an independently verified assessment.
\subsection{Short English statement}
For every integer polynomial of degree at least two, the forward orbit of any integer with infinite orbit is conjectured to contain only finitely many positive prime numbers.
\subsection{Sources}
\begin{itemize}
\item Local UnsolvedMath registry, permanent ID \(800426\), source problem code \texttt{catalogue entry}; title: \emph{Open Problems in Arithmetic Dynamics --- Conjecture 49}. Supplied status: open. Upstream revision: \texttt{d371ee917ac9647fbb06cd6fc529deec6f3dda82}.
\item Benedetto, DeMarco, Ingram, Jones et al., \emph{Current Trends and Open Problems in Arithmetic Dynamics} (2018), Conjecture 49. Supplied source URL: \url{https://arxiv.org/abs/1806.04980}. The registry attributes the problem to these authors and records the proposed year as 2018.
\item The supplied provenance records access on 2026-10-16 and extraction by \texttt{arxiv\_latex}. These are supplied provenance details; no independent source access or verification is asserted here.
\end{itemize}
\end{document}
